%% Question bank chapter file - generated by bank_tools.py build (do not edit while the bank is still being built:
%% it is regenerated from the per-test data in Question Bank\data). Every question is one bankq block:
%%   \begin{bankq}[id=..., source={...}, paper=.., qno=.., type=..., subtopic={...}, answer={...}, answer_src=key|solved,
%%                 working={...}, difficulty=1-10, diff_parts={K S I C F}, diff_note={...},
%%                 group=<passage id or empty>, also_in={duplicate copies in other tests}]
%% difficulty = min(10, K+S+I+C+F-1): K concepts 1-3, S steps 1-4, I insight 0-2, C computation 0-2, F format 0-1
%%   ...question LaTeX, options, figures (\bfig / \bfigi -> img/)...
%%   \end{bankq}
%% Passages shared by several questions are bankpassage blocks; their questions carry group=<passage id>.
\documentclass[10pt]{article}
\usepackage[a4paper,margin=1.8cm]{geometry}
\usepackage[T1]{fontenc}
\usepackage{amsmath,amssymb,mathtools}
\usepackage{graphicx}
\usepackage[export]{adjustbox}
\usepackage{xcolor,needspace,keyval}
\usepackage{enumitem}
\usepackage{booktabs,array,tabularx}
\usepackage[hidelinks]{hyperref}
\setlength{\parindent}{0pt}
\setlength{\parskip}{3pt}
\makeatletter
\define@key{bq}{id}{\def\bq@id{#1}}\define@key{bq}{source}{\def\bq@src{#1}}\define@key{bq}{paper}{\def\bq@paper{#1}}
\define@key{bq}{qno}{\def\bq@qno{#1}}\define@key{bq}{type}{\def\bq@type{#1}}\define@key{bq}{subtopic}{\def\bq@sub{#1}}
\define@key{bq}{answer}{\def\bq@ans{#1}}\define@key{bq}{answer_src}{\def\bq@asrc{#1}}\define@key{bq}{working}{\def\bq@work{#1}}
\define@key{bq}{group}{\def\bq@group{#1}}\define@key{bq}{also_in}{\def\bq@also{#1}}
\define@key{bq}{difficulty}{\def\bq@diff{#1}}\define@key{bq}{diff_parts}{\def\bq@dparts{#1}}\define@key{bq}{diff_note}{\def\bq@dnote{#1}}
\newenvironment{bankq}[1][]{\def\bq@work{}\def\bq@also{}\def\bq@group{}\def\bq@diff{?}\def\bq@dparts{}\def\bq@dnote{}\setkeys{bq}{#1}%
  \par\needspace{4\baselineskip}\medskip\noindent{\small\textbf{\bq@id}\hfill\color{gray}[\bq@src, Paper \bq@paper, Q\bq@qno]
  \bq@type\ $\cdot$ \bq@sub\ $\cdot$ \textbf{D\bq@diff}}\par\nopagebreak\relax}
  {\par{\small\color{teal}Answer: \textbf{\bq@ans}\ifx\bq@work\@empty\else\ (solved: \bq@work)\fi
  \ifx\bq@also\@empty\else\par\color{gray}Also in: \bq@also\fi}\par}
\define@key{bp}{id}{\def\bp@id{#1}}\define@key{bp}{source}{\def\bp@src{#1}}\define@key{bp}{questions}{\def\bp@qs{#1}}
\newenvironment{bankpassage}[1][]{\setkeys{bp}{#1}\par\needspace{4\baselineskip}\bigskip\noindent
  {\small\textbf{Passage \bp@id}\hfill\color{gray}[\bp@src; questions \bp@qs]}\par\nopagebreak\relax\itshape}{\par}
\makeatother
\newcommand{\bfig}[2]{\par\smallskip{\centering\includegraphics[width=#1pt,max width=0.62\linewidth,max height=0.32\textheight]{img/#2}\par}\smallskip}
\newcommand{\bfigi}[2]{\includegraphics[width=#1pt,max width=\linewidth,max height=0.32\textheight,valign=c]{img/#2}}
\newcommand{\bmissing}[1]{\textit{[figure missing in the source conversion: \detokenize{#1}]}}
\newenvironment{qoptions}{\begin{itemize}[leftmargin=1.8em,itemsep=1pt,topsep=2pt,label={}]}{\end{itemize}}
\newcommand{\opt}[2]{\item[(#1)] #2}
\providecommand{\fwstart}{}\providecommand{\fwend}{}\providecommand{\qfullsplit}{}
\newcolumntype{L}[1]{>{\raggedright\arraybackslash}p{#1}}
\newcolumntype{C}[1]{>{\centering\arraybackslash}p{#1}}
\newcolumntype{R}[1]{>{\raggedleft\arraybackslash}p{#1}}
\providecommand{\sectionhdr}[1]{}\providecommand{\subsectionhdr}[1]{}


\begin{document}
{\Large\bfseries Chemistry in Everyday Life --- Question Bank}\par\smallskip
{\small Chemistry chapter C28. Questions from all converted test papers, grouped by source series. Answers: printed key where it matches the paper, otherwise solved (answer\_src=solved, with working). Exact repeats across papers appear once, with also\_in listing the other copies.}\par\medskip\hrule

\section*{Score-I papers}
\begin{bankq}[id=SI19-P1-CHEM-S1-Q01, source={Score-I - Test 19 (20 Nov 2022)}, paper=1, qno=1, type=multi, subtopic={Drugs}, answer={A,B,C,D}, answer_src=key, difficulty=5, diff_parts={K1 S3 I1 C0 F1}, diff_note={concepts: statements on drugs soap sweeteners and polymers | steps: aspirin non narcotic analgesic; soap from fats and alkali; saccharin and alitame contain sulphur; HDPE by Ziegler Natta; check each}]
Which of the following statement(s) is/are CORRECT ?
\begin{qoptions}
\opt{A}{Aspirin is an example of non-narcotic analgesics.}
\opt{B}{Soap is formed by boiling triglyceryl stearate with caustic soda}
\opt{C}{Saccharin and alitame have S atom in their structure.}
\opt{D}{In HDPE preparation, Zeigler natta catalyest is used up}
\end{qoptions}
\end{bankq}
\begin{bankq}[id=SI19-P1-CHEM-S3-Q03, source={Score-I - Test 19 (20 Nov 2022)}, paper=1, qno=3, type=numerical, subtopic={Drugs}, answer={2}, answer_src=key, difficulty=5, diff_parts={K1 S2 I1 C1 F1}, diff_note={concepts: count bactericidal drugs | steps: penicillin ofloxacin bactericidal; erythromycin tetracycline chloramphenicol bacteriostatic; count}]
How many of the following drugs are bactericidal in nature?\\
Erythromycin, tetracycline, penicillin, ofloxacin, chloramphenicol.
\end{bankq}
\begin{bankq}[id=SI31-P2-CHEM-S1-Q16, source={Score-I - Test 31 (29 Dec 2022)}, paper=2, qno=16, type=single, subtopic={Drugs}, answer={C}, answer_src=key, difficulty=2, diff_parts={K1 S1 I1 C0 F0}, diff_note={concepts: how morphine converts into heroin | steps: acetylation of both hydroxyl groups with acetic anhydride gives diacetylmorphine}]
Morphine changes into Heroin by replacing
\begin{qoptions}
\opt{A}{all --OH by --OCH$_3$}
\opt{B}{one --OH by --OCH$_3$}
\opt{C}{all --OH by --OAc}
\opt{D}{one --OH by --OCH$_3$ and one--OH}
\end{qoptions}
\end{bankq}
\begin{bankq}[id=SI31-P2-CHEM-S2-Q10, source={Score-I - Test 31 (29 Dec 2022)}, paper=2, qno=10, type=numerical, subtopic={Drugs}, answer={10}, answer_src=solved, working={Equanil (meprobamate) is C9H18N2O4: [2x + 1] = 19 and (y + z - 2w)/2 = (18 + 4 - 4)/2 = 9, difference 10. The printed key has no value for this question.}, difficulty=4, diff_parts={K1 S1 I1 C1 F1}, diff_note={concepts: expression from the molecular formula of meprobamate | steps: Equanil is C9H18N2O4; evaluate [2x + 1] - (y + z - 2w)/2 = 19 - 9}]
Equanil is used in controlling depression and hypertension. The formula of Equanil is C$_x$H$_y$O$_z$N$_w$. Find $[2x+1]-\left(\dfrac{y+z-2w}{2}\right)$ .
\end{bankq}

\section*{Score-II papers}
\begin{bankq}[id=SII19-P1-CHEM-S1-Q17, source={Score-II - Test 19 (20 Feb 2024)}, paper=1, qno=17, type=single, subtopic={Food chemicals}, answer={B}, answer_src=key, difficulty=2, diff_parts={K1 S2 I0 C0 F0}, diff_note={concepts: match vitamins with deficiency diseases | steps: vitamin A night blindness; C scurvy; D osteomalacia or rickets; K clotting}]
Match the vitamins given in Column-I with the deficiency diseases listed in Column-II
\begin{center}
\begin{tabular}{|c|l|c|c|}
\hline
\multicolumn{2}{|c|}{\textbf{Column-I}} & \multicolumn{2}{c|}{\textbf{Column-II}}\\
\hline
(a) & Vitamin A & (p) & Osteomalacia\\
\hline
(b) & Vitamin C & (q) & Beri beri\\
\hline
(c) & Vitamin D & (r) & Scurvy\\
\hline
(d) & Vitamin B$_1$ & (s) & Xerophthalmia\\
\hline
\end{tabular}
\end{center}
\begin{qoptions}
\opt{A}{a $\rightarrow$ p; b $\rightarrow$ q; c $\rightarrow$ s; d $\rightarrow$ r}
\opt{B}{a $\rightarrow$ s; b $\rightarrow$ r; c $\rightarrow$ p; d $\rightarrow$ q}
\opt{C}{a $\rightarrow$ s; b $\rightarrow$ p; c $\rightarrow$ r; d $\rightarrow$ q}
\opt{D}{a $\rightarrow$ p; b $\rightarrow$ r; c $\rightarrow$ q; d $\rightarrow$ s}
\end{qoptions}
\end{bankq}
\begin{bankq}[id=SII36-P1-CHEM-S2-Q10, source={Score-II - Test 36 (07 Feb 2025)}, paper=1, qno=10, type=multi, subtopic={Drugs}, answer={A,C,D}, answer_src=key, difficulty=3, diff_parts={K1 S2 I0 C0 F1}, diff_note={concepts: properties of drugs and dyes | steps: penicillin chiral centres and tests; azo dye coupling conditions and use; codeine groups and DOU; aspirin synthesis from salicylic acid and acetic anhydride}]
Which are incorrectly matched?

\begin{qoptions}
\opt{A}{\textbf{Penicillin} : 3 chiral centre compound giving positive NaHCO$_3$ test and Lassign test give purple colour with sodium nitroprusside test.}
\opt{B}{\textbf{Azodye} : Used as antibiotic and formed when benzene diazonium ion reacts with phenol at pH = 8.5 - 9.5.}
\opt{C}{\textbf{Codeine} : Having phenolic OH and ether group and soluble in acid, with DOU = 11}
\opt{D}{\textbf{Aspirin} : It is used as nonnarcotic analgesic. It prevents blood clotting and it is formed when saliaylic acid reacts with methanol.}
\end{qoptions}
\end{bankq}

\section*{Score Advanced papers}
\begin{bankq}[id=SA03-P2-CHEM-S1-Q03, source={Score Advanced - Test 03 (27 Jun 2021)}, paper=2, qno=3, type=single, subtopic={Food chemicals}, answer={D}, answer_src=solved, working={vitamin D}, difficulty=1, diff_parts={K1 S1 I0 C0 F0}, diff_note={concepts: vitamins | steps: vitamin D deficiency causes rickets}]
Bone deformities in children is due to deficiency of which vitamins :
\begin{qoptions}
\opt{A}{Ascorbic acid}
\opt{B}{Thiamine}
\opt{C}{Riboflavin}
\opt{D}{Vitamin-D}
\end{qoptions}
\end{bankq}
\begin{bankq}[id=SA05-P1-CHEM-S1-Q12, source={Score Advanced - Test 05 (15 Jul 2021)}, paper=1, qno=12, type=single, subtopic={Drugs}, answer={D}, answer_src=key, difficulty=8, diff_parts={K3 S3 I2 C1 F0}, diff_note={concepts: identification of aspirin from a combustion analysis, drugs, C:H:O=7:8:1 | steps: C:H:O=7:8:1; anisole; HI gives phenol; Kolbe gives salicylic acid; acetylation}]
An organic compound A having C, H, O 0.108 gm A on combustion gives 0.308 gm CO$_2$ and 0.072 gm H$_2$O. A is insoluble in NaOH and NaHCO$_3$. A reacts with HI to give compound C and D. D reacts with neutral FeCl$_3$ to give violet color\par
\noindent D $\xrightarrow[\text{2--4 atm/H}^+]{\text{NaOH+CO}_2}$ D$_1$ $\xrightarrow{\text{CH}_3\text{COCl}}$ D$_2$\par
The product D$_2$ is
\begin{qoptions}
\opt{A}{An antiallergic}
\opt{B}{An antiseptic}
\opt{C}{An antibiotic}
\opt{D}{An antipyretic}
\end{qoptions}
\end{bankq}
\begin{bankq}[id=SA17-P2-CHEM-S2-Q09, source={Score Advanced - Test 17 (25 May 2023)}, paper=2, qno=9, type=multi, subtopic={Drugs}, answer={A,B,C,D}, answer_src=key, difficulty=5, diff_parts={K1 S2 I1 C1 F1}, diff_note={concepts: barbiturate synthesis from urea and malonic ester | steps: veronal from diethyl diethylmalonate and urea; identify P Q R S; statements about each; check}, also_in={Score Advanced - Test 35 (15 May 2024) P2 Q9}]
Choose the \textbf{CORRECT} statement(s) for given reaction scheme\par
\bfig{277.3}{SA17_p2_chem_s2_q09_1.png}

In the above reaction scheme, if [Q] is urea and [R] is veronal, then correct statement(s) is/are :-
\begin{qoptions}
\opt{A}{[R] is a derivative of barbituric acid}
\opt{B}{[R] is used as a hypnotic.}
\opt{C}{[R] contains two delocalized lone pair.}
\opt{D}{There will be three C=O groups in [R]}
\end{qoptions}
\end{bankq}

\section*{JEE Advanced (2021 Series) papers}
\begin{bankq}[id=JA05-P2-CHEM-S1-Q04, source={JEE Advanced (2021 Series) - Test 05 (11 Apr 2021)}, paper=2, qno=4, type=multi, subtopic={Drugs}, answer={A,B,C,D}, answer_src=key, difficulty=4, diff_parts={K1 S2 I1 C0 F1}, diff_note={concepts: chelating agents and metal based drugs | steps: D-penicillamine chelates copper; desferrioxamine chelates iron; EDTA calcium salt for lead poisoning; cis-platin inhibits tumour cell division}]
Which of the following statement(s) is/are CORRECT?
\begin{qoptions}
\opt{A}{D-penicillamine is used to remove excess of Cu through chelate therapy}
\opt{B}{Desferrioxamine is used to remove excess of Fe through chelate therapy}
\opt{C}{EDTA is used in the treatment of lead poisoning}
\opt{D}{cis-platin and related compound effectively inhibit the growth of tumour}
\end{qoptions}
\end{bankq}

\section*{TMAS papers}
\begin{bankq}[id=TMAS04-P1-CHEM-S1-Q06, source={TMAS - Test 04 (19 Nov 2023)}, paper=1, qno=6, type=single, subtopic={Drugs}, answer={C}, answer_src=key, difficulty=2, diff_parts={K1 S1 I0 C1 F0}, diff_note={concepts: drug structures | steps: check each structure}]
Which of the following statement is INCORRECT?
\begin{qoptions}
\opt{A}{Terfenadine has `$3$' aromatic rings in it's structure}
\opt{B}{Prontosil is coloured due to conjugated double bond system.}
\opt{C}{Chloromphenicol contains phenolic functional group.}
\opt{D}{Chloroxylenol has a phenolic functional group}
\end{qoptions}
\end{bankq}
\begin{bankq}[id=TMAS04-P1-CHEM-S2-Q09, source={TMAS - Test 04 (19 Nov 2023)}, paper=1, qno=9, type=multi, subtopic={Food chemicals}, answer={A,C,D}, answer_src=key, difficulty=3, diff_parts={K1 S1 I0 C1 F1}, diff_note={concepts: vitamins and sweeteners | steps: match each}]
\begin{center}
\begin{tabular}{|l|l|l|l|}\hline
\multicolumn{2}{|c|}{\textbf{List-I}} & \multicolumn{2}{c|}{\textbf{List-II}}\\\hline
(a) & Vitamin $\mathrm{B_2}$ & (P) & Ascorbic acid\\\hline
(b) & Vitamin C & (Q) & Artificial sweetner\\\hline
(c) & Sucralose & (R) & Thiamine\\\hline
(d) & Vitamin $\mathrm{B_1}$ & (S) & Riboflavin\\\hline
\end{tabular}
\end{center}
The incorrect match is/are?
\begin{qoptions}
\opt{A}{a$\to$P ; b$\to$Q ; c$\to$S ; d$\to$R}
\opt{B}{a$\to$S ; b$\to$P ; c$\to$Q ; d$\to$R}
\opt{C}{a$\to$Q ; b$\to$P ; c$\to$R ; d$\to$S}
\opt{D}{a$\to$R ; b$\to$Q ; c$\to$P ; d$\to$S}
\end{qoptions}
\end{bankq}
\begin{bankq}[id=TMAS04-P1-CHEM-S3-Q06, source={TMAS - Test 04 (19 Nov 2023)}, paper=1, qno=6, type=numerical, subtopic={Drugs}, answer={1.00}, answer_src=key, difficulty=3, diff_parts={K1 S1 I0 C1 F1}, diff_note={concepts: broad spectrum antibiotics | steps: classify each drug; count}]
How many of the following drugs is/are example(s) of broad spectrum antibiotic ?

Ofloxacin, Penicillin--G, Terpineol, Salvarsan
\end{bankq}

\section*{Star Review (Other) papers}
\begin{bankq}[id=STAR03-P1-CHEM-S1-Q04, source={Star Review (Other) - Test 03 (04 Dec 2023)}, paper=1, qno=4, type=single, subtopic={Drugs}, answer={C}, answer_src=key, difficulty=1, diff_parts={K1 S1 I0 C0 F0}, diff_note={concepts: chemicals acting on the stomach | steps: histamine and the antihistamines}]
A chemical which stimulates the secretion of pepsin is :
\begin{qoptions}
\opt{A}{Anti histamine}
\opt{B}{Cimetidine}
\opt{C}{Histamine}
\opt{D}{Zantac}
\end{qoptions}
\end{bankq}
\end{document}
